\nadpisOddilu{Seznam obrázků}
\addcontentsline{toc}{chapter}{Seznam obrázků}

{\color{infoBlue}
Seznam obrázků se v práci objeví za obsahem, když v práci bude pět a více obrázků. Aktualizaci seznamu provedete klávesou F9 po umístění kurzoru do textu Obsahu.
}

% Položky jako ve Wordu (Obsah1/Obsah2/Seznam obrázků: bez mezery za odstavcem)
{\setlength{\parskip}{0pt}
\makeatletter
\@starttoc{lof}
\makeatother
\par}

\chapter*{Seznam tabulek}
\addcontentsline{toc}{chapter}{Seznam tabulek}

{\color{infoBlue}
Seznam tabulek se v práci objeví zde, když v práci bude pět a více tabulek. Aktualizaci seznamu provedete klávesou F9 po umístění kurzoru do textu Obsahu.
}

% Položky jako ve Wordu (Obsah1/Obsah2/Seznam obrázků: bez mezery za odstavcem)
{\setlength{\parskip}{0pt}
\makeatletter
\@starttoc{lot}
\makeatother
\par}

\chapter*{Seznam zkratek}
\addcontentsline{toc}{chapter}{Seznam zkratek}

{\color{infoBlue}
Seznam zkratek, které nejsou běžné, se zpravidla uvádí na samostatnou stránku.
}

% Každá zkratka: \zkratka{zkratka}{význam}
\zkratka{ČR}{Česká republika}
\zkratka{ČSN}{Česká státní norma}
\zkratka{ČSÚ}{Český statistický úřad}
\zkratka{IS}{Informační systém}
\zkratka{NATO}{North Atlantic Treaty Organization (Severoatlantická aliance)}
\zkratka{OSN}{Organizace spojených národů}
\zkratka{VŠPJ}{Vysoká škola polytechnická Jihlava}

\pagebreak
